\bookchapter{A back cover with a free USB port}\label{k:cover}

Chapter~\ref{k:repo} deploys code by plugging a USB stick into the robot: no network, no laptop, which is exactly what you want in the hectic minutes before a match. On the K1 the USB port is hidden under the back cover\guess. A 3D-printable replacement cover with an opening for the port exists; this chapter prints and fits it, and checks that the robot sees a stick. Milestone: one (needed for USB deployment).

\begin{needbox}
\begin{itemize}[nosep]
  \item The model files of the replacement cover (put the link and the licence into \texttt{docs/hardware.md}), and the official CAD or robot models from Booster for comparison\guess.
  \item A 3D printer, PETG or ASA filament (PLA softens near warm electronics), the right screwdrivers, a tray for screws.
  \item A USB stick. Background: \rtr{Devices, udev Rules and Audio}.
\end{itemize}
\end{needbox}

\begin{guessbox}[Check with the rules first]
Pre-approved standard platforms are inspected to check that they are \emph{unmodified} (\rtr{Preparing for a tournament}). Whether a printed back cover counts as a modification is a question for the organisers of your competition; ask before the tournament, and keep the original cover in the team box so you can swap it back.
\end{guessbox}

\section*{Steps}
\begin{stepbox}{Step 1: compare the model with the robot}
Open the model in a slicer and in a CAD viewer next to Booster's model of the original cover. Check the positions of the screw holes, the clips, and every ventilation opening of the original: the new cover must keep them, because the computer behind it gets warm.
\end{stepbox}

\begin{stepbox}{Step 2: print a test piece}
Print only the corner with two screw holes (cut it in the slicer), $10$\,minutes instead of hours. Hold it against the robot: do the holes line up? Adjust the scale or the hole diameter in the slicer if not (printers often print holes a little too small).
\end{stepbox}

\begin{stepbox}{Step 3: print the cover}
Material PETG or ASA; walls of at least $3$ perimeters; the outside face on the print bed if it is flat (smooth surface, no supports on the visible side). Write the slicer settings into \texttt{docs/hardware.md} so the next print is the same.
\end{stepbox}

\begin{stepbox}{Step 4: swap the covers}
Robot off, battery out (or tether unplugged). Take a photo, then remove the screws into a labelled tray. Do not pull cables; if a cable is attached to the original cover (a fan, a button), note how it is plugged. Fit the new cover, tighten screws only until they hold: plastic threads strip easily.
\end{stepbox}

\begin{stepbox}{Step 5: does the robot see the stick?}
Boot, log in with \verb!ssh k1!, and watch the kernel events while you plug a stick in:
\begin{verbatim}
udevadm monitor --udev --subsystem-match=block
lsblk -o NAME,LABEL,SIZE,MOUNTPOINT
\end{verbatim}
Give the stick the label \texttt{K1DEPLOY} on the laptop (for a FAT stick: \verb!sudo fatlabel /dev/sdX1 K1DEPLOY!); \texttt{lsblk} on the robot must show that label.
\end{stepbox}

\begin{checkbox}
\begin{itemize}[nosep]
  \item The cover sits flush, all screws hold, all openings of the original are open.
  \item \texttt{udevadm monitor} prints \texttt{add} events when the stick is plugged in, and \texttt{lsblk} shows the label \texttt{K1DEPLOY}.
  \item After a $30$-minute session the cover is not deformed near the computer.
\end{itemize}
\end{checkbox}

\begin{troublebox}
\begin{description}[leftmargin=0.6cm,nosep]
  \item[No events] The port may be power-only or disabled, or the cable behind the cover is not connected. Try another stick, and compare with \verb!dmesg -w!.
  \item[The cover warps] PLA was used, or the cover blocks an air flow. Reprint in PETG or ASA and check the openings.
\end{description}
\end{troublebox}

\begin{keyidea}[What to remember]
\begin{itemize}[nosep]
  \item Compare with the original before printing; print a test piece; keep every ventilation opening.
  \item A free USB port makes deployment without network possible (Chapter~\ref{k:repo}).
  \item Ask the organisers whether the cover is allowed, and keep the original.
\end{itemize}
\end{keyidea}
